%% =============================================================
%% Section 1: Introduction
%% =============================================================

\section{Introduction}
\label{sec:intro}

Should central banks publish their interest rate forecasts? The question sits at the
heart of modern monetary policy communication. Since the Reserve Bank of New Zealand
first released an explicit rate path in 1997, a growing number of central banks have
adopted the practice of \emph{forward guidance}---communicating, in varying degrees of
detail, the expected future path of the policy rate. Yet a parallel tradition of
institutional opacity persists, and the welfare consequences of these two approaches
remain contested.

This paper develops a two-period small open economy New Keynesian model to compare
two communication regimes: \emph{explicit forward guidance}, in which the central bank
announces a forecast of the period-2 interest rate; and \emph{opacity}, in which it
remains silent about future intentions. A key feature of the model is that the central
bank's inflation aversion---its ``type''---is private information. Private agents
observe current macroeconomic conditions but do not know whether they face a hawkish
or a dovish policymaker. This uncertainty about the reaction function generates
misaligned expectations about future interest rates and, through the uncovered
interest parity (UIP) condition, misaligned exchange rate expectations in the current
period.

Forward guidance resolves this uncertainty by credibly revealing the central bank's
type. We show that this has two distinct welfare-improving channels: (i) it reduces the
variance of private-sector forecast errors about the period-2 rate, lowering expected
losses in that period; and (ii) it corrects exchange rate misalignment in the current
period, improving period-1 allocations. The exchange rate channel is specific to the
open-economy setting and amplifies the value of transparency relative to what a
closed-economy analysis would imply.

[The introduction will be extended once results are fully derived.]

\paragraph{Related literature.}
[Placeholder for literature review.]

\paragraph{Outline.}
Section~\ref{sec:model} presents the model environment, defines the two communication
regimes, and characterizes equilibrium in each period.
Section~\ref{sec:welfare} compares welfare under forward guidance and opacity.
Section~\ref{sec:extensions} discusses extensions.
Section~\ref{sec:conclusion} concludes.
